\chapter{Evaluation and Results}\label{ch:results}
% This chapter reports the results of the experimental evaluation 


% based on the study design introduced in Chapter~\ref{ch:evaluation}. It assesses how the proposed architecture maintains correct and continuous event reasoning when data becomes temporarily unavailable. 

% This chapter describes the study design used to evaluate the proposed architecture. Building on the mechanisms introduced in Chapter~\ref{ch:architecture}, it outlines the evaluation scenario, the baseline configurations, and the metrics used to assess system behaviour under temporary data unavailability.


Smart Power Grid Monitoring
Environmental Monitoring
Traffic Flow Monitoring


What you’re describing becomes a loss-of-observability disturbance model, where the main failure modes are:

DT leaves → observation disappears

DT degrades → observations become intermittent

DT becomes faulty → observations become sparse or delayed

So the physical process continues normally, but the system’s knowledge becomes incomplete.




